\section*{Capítulo \ref{ch:b3:adaptive-immunity} --- Inmunidad adaptativa y vacunación}

\begin{solution}{exo:b3:adaptive-immunity:1}
Un \omterm{def:b3:adaptive-immunity:lymphocytes}{linfocito}, una especificidad (Nossal y Lederberg: cada célula
secretaba \omterm{def:b3:adaptive-immunity:antibody}{anticuerpo} contra un solo \omterm{def:b3:adaptive-immunity:lymphocytes}{antígeno}); el receptor se fabrica
antes de encontrarse con el \omterm{def:b3:adaptive-immunity:lymphocytes}{antígeno} (Tonegawa: el gen se reordena en el
linfocito~B y no en respuesta al \omterm{def:b3:adaptive-immunity:lymphocytes}{antígeno}); el \omterm{def:b3:adaptive-immunity:lymphocytes}{antígeno} selecciona y
expande los clones que encajan en células efectoras y de memoria; y los
clones autorreactivos se eliminan durante el desarrollo.
\end{solution}

\begin{solution}{exo:b3:adaptive-immunity:2}
Dos cadenas pesadas y dos ligeras, cada una con un dominio variable en la
punta y dominios constantes por detrás; dos sitios de unión idénticos
formados por seis bucles hipervariables; y un tallo Fc que leen otras
células y el \omterm{def:b3:innate-immunity:complement}{complemento}. IgM: el primer \omterm{def:b3:adaptive-immunity:antibody}{anticuerpo}, pentámero, activa el
\omterm{def:b3:innate-immunity:complement}{complemento}. IgG: el \omterm{def:b3:adaptive-immunity:antibody}{anticuerpo} principal de la sangre y los tejidos,
\omterm{def:b3:innate-immunity:complement}{opsonización}, neutralización, atraviesa la placenta. IgA: dímero
secretado a las superficies mucosas y a la leche. IgE: unida a los
mastocitos, contra los parásitos, mediadora de la alergia. IgD: en los
linfocitos~B vírgenes, de función conocida menor.
\end{solution}

\begin{solution}{exo:b3:adaptive-immunity:3}
Clase~I: todas las células nucleadas; péptidos de proteínas citosólicas
cortados por el proteasoma; 8--10 residuos; los leen los linfocitos~T
CD8 citotóxicos. Clase~II: \omterm{def:b3:innate-immunity:innate}{células dendríticas}, \omterm{def:b3:innate-immunity:innate}{macrófagos} y
linfocitos~B; péptidos de proteínas captadas y digeridas en \omterm{def:b3:membrane-traffic:endocytosis}{endosomas};
13--25 residuos; los leen los linfocitos~T CD4 colaboradores.
\end{solution}

\begin{solution}{exo:b3:adaptive-immunity:4}
$R_{0}$ es el número medio de casos que produce un caso en una población
del todo susceptible; el umbral es $p_{c} = 1 - 1/R_{0}$. Paperas,
\qty{80}{\%}; tosferina, \qty{92}{\%}; gripe de 1918, \qty{60}{\%}.
\end{solution}

\begin{solution}{exo:b3:adaptive-immunity:5}
$50\times 20\times 5 = 5000$ cadenas pesadas y $60\times 4 = 240$
ligeras: $1.2\times 10^{6}$ combinaciones; con la \omterm{def:b3:adaptive-immunity:antibody}{diversidad de unión},
$1.2\times
10^{10}$ --- diez veces el número de linfocitos~B, de modo que el animal
expresa como mucho una décima parte de lo que podría fabricar y casi cada
célula lleva un receptor único.
\end{solution}

\begin{solution}{exo:b3:adaptive-immunity:6}
$1000\times 10^{-3} = 1$ mutación por célula y división; $20$ en veinte
divisiones. Entre $10^{5}$ células, $2\times 10^{6}$ sucesos de mutación
caen sobre $3000$ sustituciones simples posibles: cada una se muestrea
centenares de veces, y además las combinaciones dobles y triples --- el
\omterm{def:b3:adaptive-immunity:germinal-centre}{centro germinal} explora de forma exhaustiva la vecindad del receptor.
\end{solution}

\begin{solution}{exo:b3:adaptive-immunity:7}
Los linfocitos~T se seleccionan en el \omterm{def:b3:adaptive-immunity:tolerance}{timo} para reconocer péptidos sobre
los alelos MHC propios del hospedador; los alelos de otra persona tienen
surcos de forma distinta que sostienen péptidos distintos en
orientaciones distintas, de modo que ese mismo péptido vírico no se ve.
Las moléculas MHC ajenas de un injerto las reconoce a su vez una gran
parte de los linfocitos~T como ``MHC propio con un péptido extraño'', de
ahí el rechazo. El polimorfismo persiste porque una población con muchos
alelos presenta una muestra más amplia de los péptidos de cualquier
patógeno y ningún patógeno puede escapar a todo el mundo; los
heterocigotos presentan el doble de variedad y quedan favorecidos.
\end{solution}

\begin{solution}{exo:b3:adaptive-immunity:8}
(a) Sin \omterm{def:b3:adaptive-immunity:antibody}{recombinación V(D)J}: ni linfocitos~B ni T, una inmunodeficiencia
combinada grave. (b) Sin AIRE: las células tímicas no logran exponer
\omterm{def:b3:adaptive-immunity:lymphocytes}{antígenos} de los tejidos, los linfocitos~T autorreactivos escapan y hay
\omterm{def:b3:adaptive-immunity:tolerance}{autoinmunidad} contra muchos órganos. (c) Sin \omterm{def:b3:adaptive-immunity:tolerance}{FoxP3}: sin linfocitos~T
reguladores, \omterm{def:b3:adaptive-immunity:tolerance}{autoinmunidad} y alergia mortales de comienzo temprano.
(d) Sin AID: ni \omterm{def:b3:adaptive-immunity:antibody}{cambio de clase} ni hipermutación --- IgM abundante, casi
nada de IgG ni de IgA, \omterm{def:b3:adaptive-immunity:antibody}{anticuerpos} de baja afinidad e infecciones
repetidas. (e) Sin linfocitos~T CD4: sin ayuda para los linfocitos~B ni
para los \omterm{def:b3:innate-immunity:innate}{macrófagos}, respuestas de \omterm{def:b3:adaptive-immunity:antibody}{anticuerpos} pobres y susceptibilidad a
las infecciones oportunistas --- la inmunodeficiencia del sida.
\end{solution}

\begin{solution}{exo:b3:adaptive-immunity:9}
$p_{c} = 1 - 1/6 = 0.83$; cobertura $0.83/0.9 = 93\,\%$. Con una
cobertura del \qty{80}{\%} la fracción inmune es $0.72$ y
$R_{\text{ef}} = 6\times 0.28
= 1.7$. Entre los susceptibles la epidemia se propaga como si $R_{0}$
fuera $1.7$ (los inmunes se limitan a diluir los contactos), y la
relación del tamaño final $s_{\infty} - \ln s_{\infty}/1.7 = 1$ da
$s_{\infty}
\approx 0.31$: se infecta cerca del \qty{69}{\%} de los susceptibles, el
\qty{19}{\%} de la población.
\end{solution}

\begin{solution}{exo:b3:adaptive-immunity:10}
A partir de $\mathrm{d}s/\mathrm{d}t = -\beta si$ y de
$\mathrm{d}i/\mathrm{d}t
= \beta si - \gamma i$, $\mathrm{d}i/\mathrm{d}s = -1 + \gamma/(\beta s)$,
de modo que $i + s - (\ln s)/R_{0}$ se conserva; con $s = 1$ e $i = 0$ al
principio e $i = 0$ y $s = s_{\infty}$ al final, $s_{\infty} - \ln
s_{\infty}/R_{0} = 1$. Soluciones: $R_{0} = 1.5$, $s_{\infty} = 0.42$
(un \qty{58}{\%} infectado); $2$, $0.20$ (\qty{80}{\%}); $3$, $0.06$
(\qty{94}{\%}); $5$, $0.007$ (\qty{99.3}{\%}). No todo el mundo: a medida
que se agotan los susceptibles, el número reproductivo efectivo baja de
uno mientras aún quedan algunos, y la epidemia se apaga antes de
alcanzarlos.
\end{solution}

\begin{solution}{exo:b3:adaptive-immunity:11}
$1.5^{n} = 1000$: $n = \ln 1000/\ln 1.5 \approx 17$ mejoras sucesivas.
Con una mutación por célula y división y una de cada cien que mejore, hay
que cribar al menos cien células por ronda para encontrar una; en la
práctica, miles, ya que la mitad de las mutaciones destruyen la unión y
la selección es ruidosa. Diecisiete rondas de dos divisiones cada una
cuestan unos nueve días --- las dos o tres semanas observadas, dadas las
ineficiencias. Los ciclos separan la mutación (zona oscura) de la prueba
(zona clara), de modo que cada variante se juzga contra el \omterm{def:b3:adaptive-immunity:lymphocytes}{antígeno} antes
de dejarla multiplicarse, exactamente como un algoritmo genético alterna
variación y selección.
\end{solution}

\begin{solution}{exo:b3:adaptive-immunity:12}
Mujeres: varios genes inmunitarios están en el cromosoma X, algunos
(TLR7) escapan a la inactivación y se expresan desde las dos copias, y
los estrógenos favorecen la supervivencia de los linfocitos~B --- prueba:
los varones con una X de más (Klinefelter) tienen un riesgo de lupus
parecido al femenino, y así es. \omterm{def:b3:adaptive-immunity:mhc}{HLA}: los alelos asociados presentan bien
los péptidos propios pertinentes, o no logran presentarlos en el \omterm{def:b3:adaptive-immunity:tolerance}{timo} y
los clones no se eliminan --- prueba: los ratones \omterm{def:b3:genetic-engineering:transgenic}{transgénicos} para el
alelo humano desarrollan la enfermedad al recibir el \omterm{def:b3:adaptive-immunity:lymphocytes}{antígeno}. El
aumento: menos exposición microbiana temprana y \omterm{def:b3:microbiomes:symbiosis}{microbiomas} alterados
(\omterm{def:b3:bacteriology:antibiotics}{antibióticos}, dieta, parto por cesárea) dan menos linfocitos~T
reguladores y más \omterm{def:b3:innate-immunity:inflammation}{inflamación} --- pruebas: los ratones propensos a la
diabetes criados libres de gérmenes o con \omterm{def:b3:bacteriology:antibiotics}{antibióticos} la desarrollan
más, y los hijos de inmigrantes adquieren la incidencia del país nuevo en
una generación.
\end{solution}

\begin{solution}{pb:b3:adaptive-immunity:1}
\textbf{1.} $6000\times 320 = 1.9\times 10^{6}$; $\times 3\times 10^{4}
= 5.8\times 10^{10}$.
\textbf{2.} $10^{11}$ células frente a $5.8\times 10^{10}$ posibles: un
cuerpo podría albergar dos veces cada receptor, pero los clones son
desiguales y la mayoría de las secuencias nunca llegan a fabricarse, de
modo que cada cuerpo expresa una muestra grande pero incompleta --- y dos
personas comparten solo una pequeña parte de sus receptores.
\textbf{3.} $10^{11}/10^{5} = 10^{6}$ células; en un ganglio de $10^{8}$,
unas $1000$.
\textbf{4.} $10^{9}/10^{5} = 10^{4}$ precursores --- cien veces menos; el
recién nacido tiene además folículos y \omterm{def:b3:innate-immunity:innate}{células dendríticas} inmaduros, no
tiene memoria y la ayuda de los linfocitos~T es limitada, de modo que las
respuestas son lentas y pequeñas, y la IgG materna salva el intervalo.
\textbf{5.} Abarca las junturas V--D y D--J, donde se quitan y se añaden
nucleótidos al azar: su secuencia no está codificada en ningún punto del
genoma. Situado entre las dos cadenas, en el centro del sitio de unión,
es el bucle que más a menudo toca el epítopo.
\textbf{6.} Antes del \omterm{def:b3:adaptive-immunity:lymphocytes}{antígeno}, la diversidad ha de ser amplia para que
algo una lo que venga; después del \omterm{def:b3:adaptive-immunity:lymphocytes}{antígeno}, el esfuerzo se aprovecha
mejor variando los pocos receptores que ya se sabe que unen. Una búsqueda
gruesa seguida de otra fina es mucho más eficiente que cualquiera de las
dos por separado.
\textbf{7.} Siete días son $21$ divisiones: $100\times 2^{21} = 2.1\times
10^{8}$ células; el \qty{30}{\%}, $6.3\times 10^{7}$ \omterm{def:b3:adaptive-immunity:response}{células plasmáticas}.
\textbf{8.} $6.3\times 10^{7}\times 2000\times \num{86400} = 1.1\times
10^{16}$ moléculas al día; $\times 1.5\times 10^{5}\times 1.66\times
10^{-24}$ g $= \qty{2.7}{mg}$; en \qty{3}{L}, \qty{0.9}{mg/L} al día.
\textbf{9.} \qty{27}{mg} en \qty{3}{L}: \qty{9}{mg/L} $= \qty{0.009}{g/L}$,
una milésima de la IgG total --- el \omterm{def:b3:adaptive-immunity:antibody}{anticuerpo} contra un solo \omterm{def:b3:adaptive-immunity:lymphocytes}{antígeno} es
una fracción pequeña del conjunto.
\textbf{10.} Cien veces son $\log_{2}100 = 6.6$ semividas, unos $140$
días. Los años de \omterm{def:b3:adaptive-immunity:antibody}{anticuerpo} vienen de las \omterm{def:b3:adaptive-immunity:response}{células plasmáticas} longevas
de la médula, que secretan sin parar, y de las \omterm{def:b3:adaptive-immunity:response}{células de memoria}
reestimuladas al reexponerse.
\textbf{11.} Nueve divisiones: $10^{5}\times 512 = 5\times 10^{7}$
células en el día 3 --- mil veces las $5\times 10^{4}$ de la primaria en
el día 3, y una cuarta parte de lo que la primaria alcanzó solo en el día
7.
\textbf{12.} El VDJ reordenado queda junto a C$_{\mu}$, de modo que el
primer transcrito es IgM; cambiar a C$_{\gamma}$ exige ayuda de los
linfocitos~T y AID, que tardan días. Las \omterm{def:b3:adaptive-immunity:response}{células de memoria} ya han
cambiado, de modo que su progenie fabrica IgG de inmediato.
\textbf{13.} $1000\times 10^{-3} = 1$ por división; $20$ en veinte.
\textbf{14.} Una de cada cien; $10^{4}\times 0.01 = 100$ células
mejoradas por división.
\textbf{15.} $1.5^{n} = 1000$: $n \approx 17$; $17\times \qty{12}{h}
\approx 8.5$ días.
\textbf{16.} La mitad de las hijas pierden la unión: sin selección,
seguiría uniendo $0.5^{20}
\approx 10^{-6}$ del linaje al cabo de veinte divisiones. La zona clara
ha de retirar a las perdedoras en cada ciclo para que solo sigan adelante
los receptores funcionales y mejorados.
\textbf{17.} La primera dosis deja linfocitos~B de memoria de afinidad
elevada y de clase cambiada; la segunda dosis siembra con ellos \omterm{def:b3:adaptive-immunity:germinal-centre}{centros
germinales} nuevos, y la hipermutación y la selección se reanudan desde un
punto de partida más alto, de modo que los \omterm{def:b3:adaptive-immunity:antibody}{anticuerpos} que salen son aún
mejores.
\textbf{18.} Algunas personas infectadas acaban fabricando \omterm{def:b3:adaptive-immunity:antibody}{anticuerpos}
contra sitios conservados que el \omterm{def:b3:virology:virus}{virus} no puede cambiar (\omterm{def:b3:adaptive-immunity:antibody}{anticuerpos}
ampliamente neutralizantes), tras años de maduración. El \omterm{def:b3:adaptive-immunity:germinal-centre}{centro germinal}
puede dirigirse: una secuencia de inmunógenos que primero enganche a los
precursores germinales adecuados y después premie la unión al sitio
conservado podría guiar la maduración por ese camino en meses en lugar de
en años.
\textbf{19.} $p_{c} = 1 - 1/15 = 93.3\,\%$. Con una dosis: $0.933/0.93 =
100.4\,\%$ --- inalcanzable; con dos dosis: $0.933/0.97 = 96.2\,\%$.
\textbf{20.} Inmunes, $0.90\times 0.97 = 0.873$; $R_{\text{ef}} = 15\times
0.127 = 1.9$: el sarampión circula.
\textbf{21.} $60\,000\times (0.10 + 0.90\times 0.03) = 7600$ susceptibles
por cohorte de nacimiento; tras cinco años, $38\,000$ --- una reserva en
la que un caso importado encuentra entre sus contactos susceptibles
suficientes para sostener cadenas.
\textbf{22.} Entre los susceptibles, $s_{\infty} - \ln s_{\infty}/1.9 =
1$ da $s_{\infty} \approx 0.23$: alrededor del \qty{77}{\%} de ellos,
unos $29\,000$ casos.
\textbf{23.} Inmunes, $0.95\times 0.97 = 0.92$; $R_{\text{ef}} = 15\times
0.08 = 1.2$ --- todavía por encima de uno; el umbral exige un
\qty{96}{\%} con dos dosis, y el sarampión es la enfermedad que castiga
los últimos puntos porcentuales.
\textbf{24.} A los lactantes solo los protege la inmunidad de quienes los
rodean (y el \omterm{def:b3:adaptive-immunity:antibody}{anticuerpo} materno durante unos meses): si no pueden
formarse cadenas de transmisión, ningún caso les llega. Con una cobertura
del \qty{85}{\%} la fracción inmune es $0.82$ y $R_{\text{ef}} = 2.6$:
los brotes corren y los lactantes, sin vacunar y más vulnerables, se
infectan.
\textbf{25.} Repertorio de unos $6\times 10^{10}$; \qty{2.7}{mg} de IgG
al día de las \omterm{def:b3:adaptive-immunity:response}{células plasmáticas} de una respuesta; y una cobertura de
dos dosis del \qty{96}{\%}.
\end{solution}
